\chapter*{Tóm tắt Đồ án} % Hoặc bạn có thể dùng \chapter*{Abstract}

Trong bối cảnh chuyển đổi số, việc khai thác tri thức từ dữ liệu âm thanh cuộc họp gặp nhiều trở ngại do đặc thù nhiễu môi trường và sai số từ quá trình nhận dạng tiếng nói tự động (ASR). Đồ án này tập trung nghiên cứu và hiện thực hóa một quy trình xử lý đầu cuối (\textit{end-to-end pipeline}) nhằm tự động hóa việc truy xuất và tổng hợp thông tin từ các bản ghi âm hội thoại đa người nói.

Giải pháp đề xuất bao gồm một hệ thống tích hợp các mô hình Phân đoạn người nói (\textit{Speaker Diarization}) và Nhận dạng tiếng nói (ASR) làm nền tảng số hóa dữ liệu. Trọng tâm kỹ thuật nằm ở kiến trúc RAG nâng cao (\textit{Advanced RAG}) với cơ chế hiệu chỉnh ngữ nghĩa bằng Mô hình Ngôn ngữ Lớn (\textit{LLM Correction}) và chiến lược truy xuất đa bước (\textit{Multi-step Reasoning}) dựa trên phương pháp IRCoT. Hệ thống còn triển khai cơ chế tóm tắt đa tầng (\textit{Multi-level Summarization}) và lọc siêu dữ liệu (\textit{Metadata Filtering}) để tối ưu hóa tính xác thực và khả năng truy vết thông tin.

Kết quả thực nghiệm trên tập dữ liệu AMI và tập dữ liệu chuyên ngành thực tế (BIWASE) cho thấy kiến trúc Advanced RAG đạt chỉ số Hit Rate@5 lên tới 94\%, vượt trội đáng kể so với phương pháp RAG truyền thống. Hệ thống thể hiện khả năng phục hồi tri thức tốt trước các sai số ASR, đồng thời giảm mạnh kích thước dữ liệu lưu trữ với tỷ lệ nén (\textit{Compression Ratio}) đạt khoảng 0.065 mà vẫn bảo toàn được các đơn vị thông tin cốt lõi. Nghiên cứu đã chứng minh tính hiệu quả của việc kết hợp giữa tìm kiếm ngữ nghĩa và suy luận logic trong việc quản lý tri thức doanh nghiệp từ nguồn dữ liệu âm thanh phi cấu trúc.